% page 638 of the facsimile (image p-637 of the scan)
% block 144.8 x 220.4 mm ; left margin 8.0 mm
\pgc{-12.700mm}{0.000mm}{
\l{LISTO de vorti qui, pro lia teknikaleso e probiteso}
\l{multa-yara, probable adoptesos da la Akademio di Ido}
\l{---------------------------------------------------}
\l{}
\l{}
\l{\sou{chevaliero}.- I. (En\cel{1}\sou{Roma, antique})Civitano di un ek}
\l{la tri stato-klasi, inter la patrico e la plebeyo. -}
\l{II. (\sou{mezepoko}). Membro di ta insticururo milito-ten-}
\l{dencema, kun karaktero religiala, establisita che}
\l{la nobelaro feudala, qua impozis a sua membri la des-}
\l{egardo di danjero, loyaleso, la protekto gratuita}
\l{di la febli, e politegeso a la dami. - III. (\sou{nun)}}
\l{Membro di ordeno honorumala, institucita da suvereno.}
\l{IV. (\sou{nun}) En ordeno honorumala plura-klasa, membro}
\l{di la klaso minim-alta.}
\l{}
\l{\sou{doblar}. (\sou{trans.}) (\sou{Cinemo-enrejistro)}.\cel{2}(\sou{pri parol-filmo})}
\l{Tradukar lo dicata, preske vortope (remplase di la pro-}
\l{nuncatajo originala ye altra linguo), sorgante selektar}
\l{vorti qui korespondas ne nur a la senco e sentimento-}
\l{nuanci, ma anke(same-longa)a la formo lor-pronunca}
\l{di la boko di la aktoro.}
\l{}
\l{\sou{dublar}. (\sou{trans}.) Preter-vehar (generale per automobilo)}
\l{altra duktanto, pasante latere di ica lineo-dop qua}
\l{onu avancis.}
\l{}
\l{\sou{frendo}. - Kun-partoprenanto di +muvmento.}
\l{}
\l{\sou{frua}. Nemulta tempo pos la komenco di la periodo alu-}
\l{dato.}
\l{}
\l{\sou{grandoro.(matem.)}\cel{3}Quanto kontinua, qua eventuale}
\l{kreskas o deskreskas.}
\l{}
\l{\sou{identifikar}.(trans.) Komparar persono kun la deskripto}
\l{o fotografuro donata da lua pasporto od altra dokumen-}
\l{to, por divenar certena ke lu esas ya la persono quan}
\l{lu afirmas esar.}
\l{}
\l{\sou{kruzar}. (\sou{trans.}) Pasar apud (ulu), irante inversa-sinse,}
\l{quaze por atingar la loko de ube la altra persono de-}
\l{partis.}
\l{}
\l{\sou{mashelo}. (\sou{anat.}) Ensemblo ek la du parti osta di la boko}
\l{(maxilo e mandibulo) qua esas garnisita ye denti che}
\l{la vertebrozi.}
\l{}
\l{\sou{naturele}. (\sou{adverbo)}.\cel{2}Pro la naturo di la kozi.}
\l{}
\l{\sou{nexta}.\cel{2}Qua venas nemediate - spacale, tempale o rangale.}
}
